% FILE: 02_lineage_and_context.tex
% PROJECT: sources-papers
% ROLE: academic-paper-classification-and-evidence-grounding

\section{Lineage and Context}
\label{sec:sources-papers-lineage}

\subsection{2016 Language-Model Lesson}
\label{sec:sources-papers-lineage-2016}

The lineage of this project begins with a 2016 language-model experiment whose
lesson was unambiguous: local text imitation does not conserve consequence.
A model that reproduces the surface form of a published paper --- its notation,
its claim structure, its citation apparatus --- does not thereby enforce the
invariants the paper asserts. This is not a failure of the model; it is a
structural property of text generation. Text that resembles a soundness proof
is not a soundness proof.

The implication for a research program grounded in process mining literature
is direct. A system that cites van der Aalst (1998) for WF-net soundness,
Pesic and van der Aalst (2006) for DECLARE constraint templates, and Adriansyah
et al.\ (2011) for alignment-based conformance does not thereby enforce any of
those results. The citation is surface form. The enforcement is a compile-time
type constraint or a verifiable execution boundary --- and the two are
categorically distinct.

The paper corpus in \texttt{sources/papers} is the direct institutional
response to this lesson. Its purpose is not to summarize the literature but to
convert it: each paper's formal invariants become obligations in the type
system, and each paper's algorithms become graduation requirements for the
execution engine. The academic citation is the starting point; the Rust
compile-fail fixture is the terminal artifact.

\subsection{Enterprise Process Gap}
\label{sec:sources-papers-lineage-enterprise}

The enterprise process gap that motivates the broader research program is the
gap between declared process models and observable process behavior. In
enterprise deployments, process models are maintained in documentation systems,
BPMN editors, or ERP workflow configurations. Event logs accumulate in
transactional databases, message queues, and OTel trace stores. The declared
model and the observable behavior are rarely reconciled, and when they
diverge --- as the Van der Aalst Constitution requires us to treat as a
\emph{defect}, not a discrepancy --- the divergence is structurally invisible
to any system that does not perform model-vs-log comparison.

The 81-paper classification in \texttt{sources/papers} maps the academic
response to this gap across three decades of process mining research. Alpha
Miner (van der Aalst, Weijters, Maruster 2004) discovers Petri nets from event
logs. Inductive Miner (Leemans, Fahland, van der Aalst 2013) guarantees
soundness by inductive synthesis. Alignment-based conformance (Adriansyah
et al.\ 2011) quantifies the cost-optimal deviation between log behavior and
model behavior. OCEL 2.0 (van der Aalst et al.\ 2023) extends the log model
to object-centric traces to capture the multi-object reality of enterprise
processes. Each of these represents a formal response to the enterprise process
gap; the paper corpus documents what wasm4pm must implement to make those
responses machine-enforceable in production.

\subsection{Relationship to the Chatman Equation}
\label{sec:sources-papers-lineage-chatman}

The Chatman Equation, $A = \mu(O^*)$, states that autonomy is a function of
the process-intelligence operating system. For this equation to be
interpretable, $\mu$ must be a measurable process-intelligence function ---
not a declared one. The paper corpus provides the theoretical foundation for
what a measurable $\mu$ requires: a discovery algorithm that produces a typed
model from a typed event log; a conformance algorithm that produces a
$[0,1]$-bounded fitness metric; a receipt chain that binds the model, the log,
and the fitness score into a non-forgeable evidence artifact.

Specifically, the \texttt{Between01<NUM, DEN>} const-generic fraction type
documented in \texttt{PAPER\_TO\_TYPE\_LAW.md} --- which bounds all
conformance metrics in $[0,1]$ --- is the direct Rust encoding of the
measurement requirement in the Chatman Equation. A system that reports fitness
as a floating-point number without a compile-time bound is not producing a
measurement in the sense the equation requires; it is producing an assertion.
The paper corpus establishes what the type-level distinction between measurement
and assertion looks like in practice.

\subsection{Relationship to Receipts and Replay}
\label{sec:sources-papers-lineage-receipts}

Receipts and replay are the enforcement mechanism that converts paper-derived
type laws into operational evidence. The PAPER\_CANON\_RECEIPT documents the
classified paper inventory and is the first link in the lawful evidence chain.
Subsequent receipts --- PM4PY\_ORACLE\_RECEIPT, COMPAT\_LAW\_RECEIPT,
WASM4PM\_EXECUTION\_RECEIPT --- depend on the paper corpus having been
completed first.

The core\_lifecycle\_algorithms.md formalizes the Receipt-Bearing Execution
algorithm as the fourth of four core lifecycle algorithms: after Paper-to-Law
mapping, PM4Py Oracle mapping, and Admission Gate enforcement, every
board-admissible claim must pass through receipt-bearing execution that
cryptographically binds the evidence chain. The \texttt{Admitted<T, W>} type
prevents raw event laundering by requiring that any event log entering the
execution pipeline carry an admission witness --- itself grounded in the
type-law obligations derived from the paper corpus.

The adversarial-canon.md extends this to a Byzantine threat model: BLAKE3
cryptographic receipts replace case-ID correlation as the provenance mechanism,
making Sybil attacks detectable at the type level rather than the algorithmic
level. This extension is documented in the paper corpus as a research
requirement; its implementation in wasm4pm remains an open gate
(\textit{AUTHOR THESIS}: this gap between documented requirement and
implementation evidence is a primary finding of the PARTIAL status assessment).
